\documentclass{amsart}
\usepackage{amsmath,amssymb,amsthm,hyperref}
\newtheorem{theorem}{Theorem}
\newtheorem{lemma}{Lemma}
\newtheorem{proposition}{Proposition}
\theoremstyle{remark}
\newtheorem{remark}{Remark}
\begin{document}

\title{Decomposing $10$ subject to a quadratic/quartic Diophantine condition}
\maketitle

\section{Statement and conventions}

Let $n$ be a given positive number. We seek all pairs of positive reals
$(x,y)$ with
\begin{align}
x+y&=10, \label{eq:sum}\\
0<x<y, \label{eq:order}\\
\Bigl(y+\tfrac{x}{y}\Bigr)\,x&=n. \label{eq:cond}
\end{align}
We read \eqref{eq:cond} literally: the quantity $y+\dfrac{x}{y}$ is multiplied
by $x$. (The alternative parenthesisation $y+\dfrac{x}{y}\cdot x=y+\dfrac{x^{2}}{y}=n$
is treated separately in Remark~\ref{rem:alt}.)

Because $x+y=10$ and $0<x<y$, the variable $y$ is determined by $x$ via
$y=10-x$, and the order condition $x<y$ is equivalent to $x<5$. Throughout,
$x$ ranges over the open interval $I:=(0,5)$, and then $y=10-x\in(5,10)$ is
automatically positive. We will show that \eqref{eq:cond} sets up a bijection
between $x\in I$ and $n\in(0,30)$, derive the explicit dependence, and settle
the rational and integral cases.

\section{Reduction to a cubic and the governing function}

\begin{lemma}\label{lem:cubic}
For $x\in I$ (so $y=10-x>0$), condition \eqref{eq:cond} is equivalent to
\begin{equation}\label{eq:thecubic}
P_n(x):=x^{3}-19x^{2}+(n+100)x-10n=0 .
\end{equation}
Equivalently, $n=g(x)$ where
\begin{equation}\label{eq:gdef}
g(x):=\frac{x\,(x^{2}-19x+100)}{10-x}.
\end{equation}
\end{lemma}

\begin{proof}
With $y=10-x>0$, multiply \eqref{eq:cond} through by $y$ (a positive number, so
the step is reversible):
\[
\Bigl(y+\tfrac{x}{y}\Bigr)x=n
\;\Longleftrightarrow\;
xy^{2}+x^{2}=ny .
\]
Substituting $y=10-x$,
\[
x(10-x)^{2}+x^{2}=n(10-x).
\]
Expanding the left side, $x(100-20x+x^{2})+x^{2}=x^{3}-19x^{2}+100x$, so the
equation becomes
\[
x^{3}-19x^{2}+100x=n(10-x),
\]
i.e. $x^{3}-19x^{2}+(n+100)x-10n=0$, which is \eqref{eq:thecubic}. Since the
coefficient of $n$ in $P_n(x)$ is $(x-10)\neq0$ on $I$, we may solve for $n$:
\[
n=\frac{-\,(x^{3}-19x^{2}+100x)}{x-10}
 =\frac{x^{3}-19x^{2}+100x}{10-x}
 =\frac{x(x^{2}-19x+100)}{10-x}=g(x),
\]
which is \eqref{eq:gdef}. All steps are reversible for $x\in I$ because
$y=10-x>0$ there.
\end{proof}

\section{Existence, uniqueness and the admissible range of \texorpdfstring{$n$}{n}}

\begin{lemma}\label{lem:mono}
The function $g$ in \eqref{eq:gdef} is a strictly increasing bijection from
$I=(0,5)$ onto $(0,30)$, with $g(0^{+})=0$ and $g(5^{-})=30$.
\end{lemma}

\begin{proof}
On $I$ we have $10-x>0$, $x>0$, and $x^{2}-19x+100>0$ (its discriminant is
$19^{2}-4\cdot100=-39<0$, so it is positive for all real $x$). Hence
$g(x)>0$ on $I$.

For monotonicity, write
$g(x)=\dfrac{x(x^{2}-19x+100)}{10-x}=\dfrac{-x^{3}+19x^{2}-100x}{x-10}$ and
apply the quotient rule:
\[
g'(x)=\frac{-2x^{3}+49x^{2}-380x+1000}{(x-10)^{2}}
     =\frac{-(2x-25)\,(x^{2}-12x+40)}{(x-10)^{2}} .
\]
(The factorisation is checked by expansion:
$-(2x-25)(x^{2}-12x+40)=-2x^{3}+49x^{2}-380x+1000$.) For $x\in I=(0,5)$:
\begin{itemize}
\item $2x-25<0$, hence $-(2x-25)>0$;
\item $x^{2}-12x+40$ has discriminant $144-160=-16<0$, so it is positive for
      all real $x$;
\item $(x-10)^{2}>0$.
\end{itemize}
Therefore $g'(x)>0$ for every $x\in I$, so $g$ is strictly increasing on $I$.

Finally,
\[
g(0^{+})=\frac{0\cdot100}{10}=0,
\qquad
g(5^{-})=\frac{5\,(25-95+100)}{10-5}=\frac{5\cdot30}{5}=30 .
\]
By continuity and strict monotonicity, $g$ maps $I$ bijectively onto the open
interval $(0,30)$.
\end{proof}

\begin{theorem}[Existence and uniqueness]\label{thm:main}
Let $n>0$.
\begin{enumerate}
\item If $n\in(0,30)$, the system \eqref{eq:sum}--\eqref{eq:cond} has exactly
one solution $(x,y)$. Its smaller part $x$ is the unique root of the cubic
\eqref{eq:thecubic} lying in $(0,5)$, and $y=10-x$.
\item If $n\ge 30$, the system has no solution.
\end{enumerate}
\end{theorem}

\begin{proof}
By Lemma~\ref{lem:cubic}, admissible solutions correspond exactly to
$x\in I$ with $g(x)=n$, together with $y=10-x$. By Lemma~\ref{lem:mono},
$g\colon I\to(0,30)$ is a strictly increasing bijection. Hence:
if $n\in(0,30)$ there is exactly one $x=g^{-1}(n)\in I$, giving exactly one
admissible pair; this $x$ is by Lemma~\ref{lem:cubic} a root of $P_n$, and it is
the only root in $(0,5)$ because $g$ is injective there. If $n\notin(0,30)$,
$n$ is not in the range of $g$ on $I$, so no admissible $x$ exists.
\end{proof}

\begin{remark}
The borderline value $n=30$ corresponds to $x=y=5$, which violates the strict
inequality $x<y$; that is why $n=30$ is excluded. As $n\downarrow0$ the smaller
part $x\downarrow0$.
\end{remark}

\section{Explicit formula for \texorpdfstring{$x$}{x} and \texorpdfstring{$y$}{y}}

Put $x=t+\tfrac{19}{3}$ to depress \eqref{eq:thecubic} to $t^{3}+pt+q=0$.
A direct substitution and expansion give
\begin{equation}\label{eq:pq}
p=n-\frac{61}{3},
\qquad
q=\frac{3382-99\,n}{27}\;\Bigl(=\tfrac{3382}{27}-\tfrac{11}{3}\,n\Bigr).
\end{equation}
The radicand of Cardano's formula is
\[
R:=\frac{q^{2}}{4}+\frac{p^{3}}{27}.
\]
A computation (equivalently, $-108R$ equals the cubic discriminant
$-4p^{3}-27q^{2}=-4n^{3}-119n^{2}+19840n-390000$) shows
\[
108\,R=4n^{3}+119n^{2}-19840n+390000 .
\]
Its real roots are at $n=-\tfrac{375}{4}=-93.75$ and a complex pair
$32\pm 4i$; consequently the cubic $108R$ in $n$ has a single real zero at
$n=-93.75$, and since its leading coefficient is positive and
$108R\big|_{n=0}=390000>0$, we get $R>0$ for every $n\in[0,30]$ (in particular
for all $n\in(0,30)$). Hence on the admissible range the cubic $P_n$ has
\emph{one} real root and a conjugate pair of complex roots.

\paragraph{The real root via real cube roots.}
Because $R>0$, the two numbers
\[
u_+:=-\frac{q}{2}+\sqrt{R},\qquad u_-:=-\frac{q}{2}-\sqrt{R}
\]
are \emph{real} (and $u_+\,u_-=\tfrac{q^2}{4}-R=-\tfrac{p^3}{27}$ may have either
sign). Let $\sqrt[3]{\,\cdot\,}$ denote the unique \emph{real} cube root of a
real number (defined for all reals, with $\sqrt[3]{-a}=-\sqrt[3]{a}$). Set
\[
A:=\sqrt[3]{u_+},\qquad B:=\sqrt[3]{u_-},\qquad t_0:=A+B .
\]
We check directly that $t_0$ solves the depressed cubic $t^3+pt+q=0$. Indeed,
\[
A^3+B^3=u_++u_-=-q,\qquad
AB=\sqrt[3]{u_+u_-}=\sqrt[3]{-\tfrac{p^3}{27}}=-\frac{p}{3},
\]
the second identity holding for the real cube roots since
$\bigl(-\tfrac{p}{3}\bigr)^3=-\tfrac{p^3}{27}=u_+u_-$ and the real cube root is
the unique real solution of $z^3=u_+u_-$. Therefore
\[
t_0^{3}=(A+B)^3=A^3+B^3+3AB(A+B)=-q+3\Bigl(-\tfrac{p}{3}\Bigr)t_0=-q-p\,t_0,
\]
i.e. $t_0^{3}+p\,t_0+q=0$. Thus $t_0$ is a real root of the depressed cubic, and
$x_0:=t_0+\tfrac{19}{3}$ is a real root of $P_n$. Since $P_n$ has exactly one
real root (as $R>0$), $x_0$ is \emph{the} unique real root, which by
Theorem~\ref{thm:main} is the required smaller part $x\in(0,5)$. We conclude
\begin{equation}\label{eq:cardano}
\boxed{\,x=\frac{19}{3}
   +\sqrt[3]{-\frac{q}{2}+\sqrt{R}}
   +\sqrt[3]{-\frac{q}{2}-\sqrt{R}}\,,
\qquad y=10-x\,,}
\end{equation}
with $p,q$ as in \eqref{eq:pq}, $R=\frac{q^2}{4}+\frac{p^3}{27}>0$, and the
\emph{real} cube roots taken throughout. (For instance, at $n=15$ one has
$p=-\tfrac{16}{3}$, $q=\tfrac{1897}{27}$, numerically
$R=1228.47\ldots>0$, $u_+=-0.0801\ldots$, $u_-=-70.179\ldots$, hence
$x=1.77753\ldots$, $y=8.22246\ldots$, and direct substitution gives
$(y+x/y)x=15$. Note that $u_-<0$ here, so the use of \emph{real} cube roots is
essential: the principal complex cube root would not return the real root.)

\begin{remark}
The appearance of the square root of a polynomial (the radicand $R$ is, up to a
constant, a cubic in $n$ whose relevant structure is governed by the quadratic
factor producing the complex pair $32\pm4i$) under a cube root is exactly the
``quadratic/quartic'' radical flavour referenced in the title: there is no
radical-free closed form for $x$ as a function of $n$, in contrast to the
cleanly quadratic alternative reading of Remark~\ref{rem:alt}.
\end{remark}

\section{Rational and integral solvability}

\begin{proposition}[Rational solvability]\label{prop:rat}
Fix $n\in(0,30)$ and let $(x,y)$ be the unique solution of
Theorem~\ref{thm:main}.
\begin{enumerate}
\item $x$ is rational if and only if $y=10-x$ is rational.
\item If $x$ is rational then $n=g(x)$ is rational; thus rational
decompositions occur only for rational $n$.
\item The set of values $n$ for which $x$ (equivalently $y$) is rational is
exactly
\[
\mathcal{R}=\bigl\{\,g(r)\;:\; r\in\mathbb{Q}\cap(0,5)\,\bigr\},
\]
where $g$ is given by \eqref{eq:gdef}. For such $n$, the rational $x$ is the
unique rational root of $P_n$ in $(0,5)$, namely $x=r$, and $y=10-r$.
\end{enumerate}
\end{proposition}

\begin{proof}
(1) is immediate since $y=10-x$ and $10\in\mathbb{Q}$.

(2) If $x\in\mathbb{Q}$ then $n=g(x)$ is a rational function of $x$ with
rational coefficients (by \eqref{eq:gdef}), hence rational.

(3) If $n=g(r)$ with $r\in\mathbb{Q}\cap(0,5)$, then $r$ is a root of $P_n$
in $(0,5)$ by Lemma~\ref{lem:cubic}; by uniqueness (Theorem~\ref{thm:main})
the solution is $x=r\in\mathbb{Q}$, $y=10-r$. Conversely, if for some
$n\in(0,30)$ the solution $x$ is rational, then $x\in\mathbb{Q}\cap(0,5)$ and
$n=g(x)\in\mathcal R$ by Lemma~\ref{lem:cubic}. Thus $n$ admits a rational
decomposition iff $n\in\mathcal R$.
\end{proof}

\begin{proposition}[Integral solvability]\label{prop:int}
There is \emph{no} value of $n$ for which both parts $x,y$ are positive
integers satisfying \eqref{eq:order} and \eqref{eq:cond}. The only candidate
integer smaller parts are $x\in\{1,2,3,4\}$ (forced by $0<x<5$), and they give
\[
\begin{array}{c|cccc}
x & 1 & 2 & 3 & 4\\\hline
y=10-x & 9 & 8 & 7 & 6\\[2pt]
n=g(x) & \dfrac{82}{9} & \dfrac{33}{2} & \dfrac{156}{7} & \dfrac{80}{3}
\end{array}
\]
none of which is an integer. Hence no integer $n$ admits an integral
decomposition, and each of the four admissible integer pairs $(x,y)$ produces a
non-integer rational value of $n$.
\end{proposition}

\begin{proof}
An integral decomposition requires $x,y\in\mathbb{Z}_{>0}$ with $x+y=10$ and
$x<y$, hence $x\in\{1,2,3,4\}$ and $y=10-x$. Evaluating $g$ from
\eqref{eq:gdef}:
\[
g(1)=\frac{1(1-19+100)}{9}=\frac{82}{9},\qquad
g(2)=\frac{2(4-38+100)}{8}=\frac{33}{2},
\]
\[
g(3)=\frac{3(9-57+100)}{7}=\frac{156}{7},\qquad
g(4)=\frac{4(16-76+100)}{6}=\frac{80}{3}.
\]
Each value is a non-integer rational, so no integer $n$ corresponds to an
integral $(x,y)$. (These four pairs do, however, furnish rational solutions of
Proposition~\ref{prop:rat} for the displayed rational $n$.)
\end{proof}

\section{The alternative reading}\label{sec:alt}

\begin{remark}\label{rem:alt}
If condition \eqref{eq:cond} is instead read as
$y+\dfrac{x}{y}\cdot x=y+\dfrac{x^{2}}{y}=n$, then, multiplying by $y=10-x>0$,
\[
(10-x)^{2}+x^{2}=n(10-x)
\;\Longleftrightarrow\;
2x^{2}+(n-20)x+(100-10n)=0 .
\]
This is a genuine quadratic in $x$ with discriminant
\[
D=(n-20)^{2}-8(100-10n)=n^{2}+40n-400 .
\]
Its two roots in $n$ are $-20\pm20\sqrt2$, so $D\ge0$ for $n>0$ exactly when
$n\ge 20(\sqrt2-1)=8.284\ldots$. When $D\ge0$ the two roots in $x$ are
\[
x=\frac{20-n\pm\sqrt{n^{2}+40n-400}}{4},
\]
and the \emph{smaller} candidate part is
$x_{-}=\dfrac{20-n-\sqrt{n^{2}+40n-400}}{4}$, $y=10-x_{-}$. However, one must
also impose the positivity and ordering constraints $0<x_-<5$. A short check
shows $x_-=0$ at $n=10$ and $x_-<0$ for $n>10$; hence in this alternative
reading admissible solutions exist only for
$n\in\bigl[20(\sqrt2-1),\,10\bigr)$ (for which $0<x_-<5<y$). Within that range
$x_-$ is rational iff $n$ is rational and $n^{2}+40n-400$ is a perfect rational
square, giving the clean quadratic Diophantine characterisation suggested by the
title.
We record this reading only for completeness; the literal reading of
Sections~2--5 is the one solved in full above.
\end{remark}

\end{document}
